\section{Contested legacy and final synthesis}\label{sec:peter-synthesis}

Peter's historical significance can scarcely be separated from arguments
about authority. Catholic Christians see in him the rock and shepherd whose
ministry continues in the Bishop of Rome. Orthodox Christians honor his
primacy among the apostles while disputing Catholic accounts of universal
papal jurisdiction. Protestant traditions venerate the apostle and read the
rock, keys, and succession through diverse ecclesiologies. Historical study
cannot settle those theological divisions by counting Gospel mentions; it can
prevent all sides from arguing with caricatures.

\subsection{Power under judgment}

Peter's name has authorized unity, missionary courage, doctrinal continuity,
and pastoral service. It has also been invoked to shield officeholders from
criticism or to identify a contingent Roman policy with the apostle himself.
The canonical Peter resists that misuse. He is rebuked by Jesus, receives a
commission after denial, speaks in council rather than eliminating council,
and is publicly opposed by Paul. First Peter tells elders not to lord it over
those entrusted to them. Petrine authority is credible in the pattern of the
crucified Shepherd, not as escape from accountability.

The deaths of Ananias and Sapphira and the rebuke of Simon Magus have likewise
supported severe images of ecclesial power. Acts presents divine judgment and
the non-commercial character of grace. It does not authorize manipulation,
summary violence, or purchase of office. Peter's own sword is put away. A
hagiography that turns these episodes into celebrations of domination has
missed their canonical frame.

\subsection{Peter, Israel, and the nations}

Peter's Jewish identity is not an immature stage he sheds. His confession,
Scriptural preaching, Temple prayer, dietary practice, mission to the
circumcised, and argument about Gentile fellowship all belong inside the
first-century Jewish world. Acts' claim that Gentiles need not bear the full
yoke at issue cannot responsibly be paraphrased as contempt for Torah or as
evidence that God rejected Israel. Paul calls Peter's table withdrawal a
failure precisely because the shared life created in Israel's Messiah was at
stake, not because Jewish practice was intrinsically hypocrisy.

Christian preaching has at times extracted accusations from Acts and the
Passion narratives to assign collective guilt to Jews. That reception is
theologically and historically indefensible. The Second Vatican Council's
\emph{Nostra aetate} rejects the charge against all Jews then living or Jews
today and teaches that God's gifts and calling are irrevocable. Reading Peter
within Israel is therefore not a modern courtesy added to the sources; it is
necessary to resist a contradiction of the Church's own teaching.

\subsection{A concord that includes disagreement}

The joined cult of Peter and Paul has sometimes been told as though their
agreement erased conflict. Their concord is more demanding. Paul sought
Cephas, received fellowship, opposed him publicly, and continued to recognize
his apostolate. Roman memory joined two missions rather than flattening them.
The two biographies in this collection are accordingly separate: Paul is not
an appendix to Peter's office, and Peter is not a foil for Paul's gospel. Their
shared feast is strongest when the distinct evidence remains audible.

\subsection{The shape of Peter's sanctity}

The historically recoverable Peter is fragmentary but not vague. He was a
married Jewish fisherman from the Galilean lake, one of Jesus' earliest
followers, and the most frequently foregrounded member of the Twelve. He
confessed Jesus as Messiah and resisted the road by which the Messiah would
reign. He denied Jesus and became a primary Easter witness. He led public
proclamation in Jerusalem, took part in mission beyond it, recognized the
Spirit's gift to Gentiles, and faltered over the table fellowship that gift
required. He was known personally to Paul and remembered in early, independent
streams. A converging Roman tradition makes his martyrdom under Nero probable,
while the exact execution scene remains beyond recovery. By the second
century, Christians marked a site at the Vatican as his memorial.

The hagiographic Peter is not a substitute for this man but the Church's
interpretation of him. The keys, boat, rooster, tears, chains, and inverted
cross arrange his life around mercy and service. Some details are canonical,
some early received, and some late. At their best they confess that stability
in the Church is not human self-sufficiency. The rock is a forgiven witness
whose strength is prayer received from Christ.

Unresolved questions remain substantial: the composition of the Petrine
letters; the route and duration of his ministries after Jerusalem; the exact
form of leadership in Antioch and Rome; the date and manner of death; the
history of the Vatican remains. None should be hidden. Nor do they dissolve
the central convergence: Peter's life connects the ministry of Jesus, the
Easter proclamation, the painful expansion toward the nations, and the Roman
martyr memory in a way no later legend could have created from nothing.

\begin{studybox}{Rock and weakness}
Peter is called rock without ceasing to be Simon. His biography makes the
paradox concrete: a real man's weakness is neither denied nor enthroned, but
taken into a vocation of strengthening others. For Christian memory, that is
why his authority and his martyrdom belong together.
\end{studybox}
